\documentclass[11pt]{article}
\usepackage[T1]{fontenc}
\usepackage{lmodern}
\usepackage{amsmath,amssymb,amsthm,mathtools}
\usepackage[margin=1in]{geometry}
\usepackage{microtype,xurl}
\usepackage[colorlinks=true,linkcolor=blue,citecolor=blue,urlcolor=blue]{hyperref}
\newtheorem{theorem}{Theorem}[section]
\newtheorem{lemma}[theorem]{Lemma}
\newtheorem{proposition}[theorem]{Proposition}
\theoremstyle{remark}
\newtheorem{remark}[theorem]{Remark}
\newcommand{\E}{\mathbb E}
\newcommand{\R}{\mathbb R}
\newcommand{\F}{\mathcal F}
\newcommand{\Acal}{\mathcal A}
\newcommand{\norm}[1]{\lVert#1\rVert}
\newcommand{\abs}[1]{\lvert#1\rvert}
\newcommand{\ind}{\mathbf 1}
\DeclareMathOperator{\TV}{TV}
\hypersetup{pdftitle={Strong convergence of the fractional WIS Euler integrator for Lipschitz drifts},
 pdfauthor={Kaiwen Jin},pdfsubject={Terminal strong convergence for a fixed Wick-Ito-Skorohod Euler scheme}}
\title{Strong convergence of the fractional WIS Euler integrator for Lipschitz drifts}
\author{Kaiwen Jin\\ \small\texttt{yui.ui.siki@gmail.com}}
\date{October 13, 2026}
\begin{document}
\maketitle

\begin{abstract}
We prove a terminal strong convergence bound for the GfBm\_EM Euler
integrator of a scalar stochastic differential equation driven by
multiplicative fractional Brownian noise in the Wick--It\^o--Skorohod
sense. For every Hurst parameter \(H\in(0,1)\) and every autonomous
globally Lipschitz drift, the root mean-square error is bounded by
\(C h^{\min\{H+1/2,1\}}\), using exact samples of the same noise path.
The drift need not have higher derivatives. Initial values may be zero,
and solutions may cross zero. Conditional Gaussian averaging controls the
signed quadrature error at each time scale. Summing those scales gives
the fractional rate without an exponent loss. At the Brownian
endpoint, time reversal and short intervals chosen by the distance
from zero yield order one with polynomial localization constants.
\end{abstract}

\section{Introduction}

Fractional Brownian motion provides a Gaussian driving noise with
correlated increments, but multiplicative equations require an
interpretation of the stochastic integral. The Wick--It\^o--Skorohod
(WIS) interpretation permits all Hurst parameters \(H\in(0,1)\) and
leads to a useful exact integrating-factor representation.
Erdo\u{g}an, Lord and Schieven \cite{ELS} used this representation to
define and implement the GfBm\_EM method for scalar equations with
linear multiplicative noise. Their method applies ordinary Euler
time stepping to a transformed random ordinary differential equation.
It is exact for the homogeneous linear equation.

For autonomous globally Lipschitz drifts, Theorem 3.6 of \cite{ELS}
proves a root mean-square terminal error of order \(h^H\).
The same paper conjectures the sharper rate
\(h^{\min\{H+1/2,1\}}\) in equation (1.3).
Its Section 6 explains the difficulty: local noise defects are small,
but their random coefficients prevent the cross terms from being
discarded as martingale orthogonality. Summing absolute cell errors
therefore loses the improvement seen in the signed time integral.

The present paper proves the conjectured autonomous estimate for the
stated scheme and data class. More precisely, for every fixed
\(T>0\), \(H\in(0,1)\), real coefficients and initial state, and
autonomous globally Lipschitz drift of at most linear growth, the
terminal approximation using the exact noise samples has error
\[
 \bigl(\E\abs{X_T-X_n}^2\bigr)^{1/2}
 \le C(T/n)^{\min\{H+1/2,1\}} .
\]
The constant is independent of \(n\). We require no differentiability
of the drift and include zero initial states and crossings of zero.
The exact recurrence and full statement are given in
Section~\ref{sec:setup}.

Related quadrature gains occur in several neighboring settings.
Butkovsky, Dareiotis and Gerencs\'er \cite{BDG} obtain nearly
\(\min\{1/2+\vartheta H,1\}\) strong rates for additive fractional
equations under suitable \(\vartheta\)-H\"older drift conditions,
using stochastic sewing. For additive equations with unbounded
Lipschitz drift, Moritoki \cite[Theorem 1.1]{Moritoki} proves order
one when \(H>1/2\) and order \(H+1/2-\varepsilon\), for every
\(\varepsilon>0\), when \(H\le1/2\).
Altmeyer \cite[Theorems 4--5]{Altmeyer} studies quadrature of
occupation functionals of unperturbed fractional Brownian motion
under Sobolev regularity. Kloeden and Rosa \cite{KR} establish
order-one Euler convergence for semimartingale-driven random
ordinary differential equations and treat a linear fractional-noise
example in Appendix A. These results identify useful averaging
mechanisms, but do not directly estimate the present quadrature with
its random Euler state and exponential integrating factor.

Two features of the transformed equation are decisive here.
Its Lipschitz constant in the state is deterministic, while the
unpredictable part of that state over an interval of length
\(\ell\) is only of order \(\ell^{1+H}\).
Conditioning in the past smooths the Gaussian noise variable.
The resulting local bound retains both
\(\ell h^H\) and \(h\ell^H\); a dyadic martingale sum then gives
\(h^{H+1/2}\) for \(H<1/2\) and \(h\) for \(H>1/2\).
Keeping the dependence on the interval length avoids replacing the
target exponent by a nearby exponent.

The Brownian endpoint is critical for that sum. We treat it separately.
Forward and backward It\^o formulas give a maximal order-one
quadrature bound for a Lipschitz function of Brownian motion.
After separating the value of the drift at zero, the remaining
noise coefficient vanishes with the transformed state.
Short grid intervals of length proportional to the square of its
distance from zero keep a local Girsanov density uniformly controlled.
Their contributions and number have summable powers of that distance;
the integrating-factor localization grows only polynomially.

Section~\ref{sec:setup} gives the exact terminal representation and
common estimates. Section~\ref{sec:fractional} proves the conditional
fractional quadrature bound. Section~\ref{sec:brownian} treats the
Brownian endpoint and completes the terminal error estimate.


\section{Terminal representation and common estimates}\label{sec:setup}

Fix \(T>0\), \(H\in(0,1)\), \(\alpha,\beta,x_0\in\R\), and an autonomous
globally Lipschitz function \(a:\R\to\R\) of at most linear growth.
Let \(B^H\) be centered fractional Brownian motion, normalized by
\[
 \E[B^H_tB^H_s]=\frac12(t^{2H}+s^{2H}-\abs{t-s}^{2H}).
\]
The solution \(X\) of
\begin{equation}\label{eq:sde}
 dX_t=(\alpha X_t+a(X_t))\,dt+\beta X_t\,dB^H_t,\qquad X_0=x_0,
\end{equation}
uses the Wick--It\^o--Skorohod (WIS) interpretation of
\cite{ELS}. Define, for this fixed terminal time,
\begin{equation}\label{eq:factor}
 c(t)=-\alpha t+\frac{\beta^2}{2}\bigl(T^{2H}-(T-t)^{2H}\bigr),
 \qquad J_t=\exp(c(t)-\beta B^H_t).
\end{equation}
Theorem 3.2 and equations (3.6)--(3.9) of \cite{ELS} give the exact
representation
\begin{equation}\label{eq:exact-ode}
 X_T=J_T^{-1}Z_T,\qquad
 Z_t=x_0+\int_0^t F(c(s),B^H_s,Z_s)\,ds,\qquad
 F(c,b,z)=e^{c-\beta b}a(e^{-c+\beta b}z).
\end{equation}
Thus the WIS integral is accounted for by the terminal factor and the
deterministic correction in \(c\). All subsequent time integrals involving
\(Z\) are ordinary Lebesgue integrals on the same noise path.

For \(n\ge1\), put \(h=T/n\), \(t_i=ih\), and
\(\eta(t)=h\lfloor t/h\rfloor\), with endpoint values immaterial in integrals.
The prescribed GfBm\_EM method is
\begin{equation}\label{eq:scheme}
 Z^n_0=x_0,\qquad
 Z^n_{t_{i+1}}=Z^n_{t_i}+hF(c(t_i),B^H_{t_i},Z^n_{t_i}),\qquad
 X_n=J_T^{-1}Z^n_T .
\end{equation}
It uses exact samples from the path appearing in \eqref{eq:exact-ode}.
We use the continuous interpolation
\begin{equation}\label{eq:interpolation}
 Z^n_t=x_0+\int_0^t F(c(\eta(s)),B^H_{\eta(s)},Z^n_{\eta(s)})\,ds.
\end{equation}

\begin{theorem}\label{thm:main}
For every fixed data set just specified, there is a finite constant \(C\),
independent of the integer \(n\ge1\), such that the scheme
\eqref{eq:scheme} satisfies
\begin{equation}\label{eq:main-rate}
 \bigl(\E\abs{X_T-X_n}^2\bigr)^{1/2}
 \le C h^{\min\{H+1/2,1\}}.
\end{equation}
\end{theorem}

The constant may depend on \(T,H,\alpha,\beta,x_0\) and \(a\).
No derivative assumption on \(a\), sign restriction on \(x_0\), or
restriction on \(a(0)\) is imposed. The estimate concerns the terminal
value and exact noise samples. Its proof appears at the end of
Section~\ref{sec:brownian}, after the two noise regimes have been treated.

Write \(\norm{Y}_p=(\E\abs{Y}^p)^{1/p}\).
Constants \(C_p\) below may depend on the fixed problem data and on \(p\);
they are independent of \(n\) and of time variables in \([0,T]\).
Let \(L\) be a Lipschitz constant for \(a\), \(a_0=a(0)\), and \(a_1=a-a_0\).
Then
\begin{equation}\label{eq:basic-F}
 \abs{F(c,b,z)}\le L\abs z+\abs{a_0}e^{c-\beta b},\qquad
 \abs{F(c,b,z)-F(c,b,z')}\le L\abs{z-z'}.
\end{equation}
The second bound is deterministic despite the random integrating factor.

The dependence on the noise can also be estimated using only \(L\).
For \(\sigma\in\{-1,1\}\), define
\begin{equation}\label{eq:ratio}
 b_\sigma(q)=\frac{a_1(\sigma e^q)}{\sigma e^q}.
\end{equation}
We have \(\abs{b_\sigma}\le L\) and
\([b_\sigma]_{\mathrm{Lip}}\le2L\). Indeed, for \(q>r\),
with \(x=\sigma e^q\), \(y=\sigma e^r\),
\[
 \left|\frac{a_1(x)}x-\frac{a_1(y)}y\right|
 \le 2L\frac{\abs{x-y}}{\abs x}
 =2L(1-e^{r-q})\le2L(q-r).
\]
For \(z\ne0\), the part of \(F\) corresponding to \(a_1\) is
\begin{equation}\label{eq:F1}
 F_1(c,b,z)=z\,b_{\operatorname{sign}z}
                  (\log\abs z-c+\beta b),
\end{equation}
and \(F_1(c,b,0)=0\). Consequently,
\begin{align}
 \abs{F_1(c,b,z)}&\le L\abs z,\label{eq:F1-bounds}\\
 \abs{F_1(c,b,z)-F_1(c,b',z)}
       &\le2L\abs\beta\,\abs z\,\abs{b-b'},\nonumber\\
 \abs{F_1(c,b,z)-F_1(c',b,z)}
       &\le2L\abs z\,\abs{c-c'}.\nonumber
\end{align}
The \(a_0\) part is \(a_0e^{c-\beta b}\).
The inequality
\(\abs{e^u-e^v}\le(e^u+e^v)\abs{u-v}\), together with
\eqref{eq:F1-bounds}, gives all noise and \(c\) increments used below.
For example, if \(z\) has bounded moments of the required orders, then
H\"older's inequality and Gaussian exponential moments give
\begin{equation}\label{eq:F-noise}
 \norm{F(c,B^H_r,z)-F(c,B^H_u,z)}_p
 \le C_p\abs{r-u}^H
\end{equation}
for \(c\) in a fixed bounded interval. Independence of \(z\) and the noise
is not required.

We will smooth the noise variable, rather than the original drift.
Let \(P_v f\) be convolution with the centered Gaussian density of
variance \(v>0\), and put \(f_{c,z}(b)=F(c,b,z)\).

\begin{lemma}\label{lem:heat}
For bounded \(c\) and \(0<v\le V<\infty\),
\begin{align}
 \abs{\partial_bP_vf_{c,z}(b)}
    &\le C(\abs z+e^{-\beta b}),\label{eq:heat-weighted}\\
 \abs{\partial_vP_vf_{c,z}(b)}
    &\le Cv^{-1/2}(\abs z+e^{-\beta b}).\nonumber
\end{align}
For \(g=f_{c,z}-f_{c,z'}\),
\begin{equation}\label{eq:heat-mixed}
 \norm{\partial_bP_vg}_\infty\le CL\abs{z-z'}v^{-1/2},
 \qquad
 \norm{\partial_vP_vg}_\infty\le CL\abs{z-z'}v^{-1}.
\end{equation}
Here the constants are independent of \(z,z',b,v\) in the stated ranges.
\end{lemma}
\begin{proof}
The increment estimates above bound \(\abs{f_{c,z}(b+y)-f_{c,z}(b)}\)
by \(C\abs y(\abs z+e^{-\beta b}(1+e^{\abs\beta\abs y}))\).
Both derivatives of the Gaussian kernel have integral zero. Subtract
\(f_{c,z}(b)\) inside their integrals. For the first derivative,
the integral of \(\abs y\) times its absolute value is bounded by a
constant; for the second it is bounded by \(Cv^{-1/2}\).
The exponential weight obeys the same bounds for \(v\le V\), by Gaussian
scaling and exponential moments. Since
\(\partial_vP_v=\tfrac12\partial_{bb}P_v\), this proves
\eqref{eq:heat-weighted}. Finally,
\(\norm g_\infty\le L\abs{z-z'}\) by \eqref{eq:basic-F}.
The \(L^1\) norms of the first spatial and variance derivatives of the
Gaussian kernel are \(Cv^{-1/2}\) and \(Cv^{-1}\), respectively, which
proves \eqref{eq:heat-mixed}. All differentiated integrals converge
absolutely, including the exponential weights.
\end{proof}

\begin{lemma}\label{lem:moments}
For every finite \(p\ge1\),
\begin{equation}\label{eq:moments}
 \sup_{n,t\le T}\norm{Z^n_t}_p+\sup_{t\le T}\norm{Z_t}_p\le C_p,
 \qquad
 \norm{Z^n_t-Z^n_s}_p+\norm{Z_t-Z_s}_p\le C_p\abs{t-s}.
\end{equation}
Also \(J_t\) and \(J_t^{-1}\) have uniformly bounded moments of every
finite order, and
\begin{equation}\label{eq:c-var}
 \TV(c;[0,T])\le\abs\alpha T+\frac{\beta^2}{2}T^{2H}.
\end{equation}
\end{lemma}
\begin{proof}
The assertions for \(J\) follow from its Gaussian exponential formula.
Equation \eqref{eq:basic-F} and Minkowski's inequality give the scalar
Gronwall bounds for \(Z\) and the discrete Gronwall bounds for the
grid values of \(Z^n\). Equation \eqref{eq:interpolation} gives the same
bound between grid points. Integrating the bounded \(L^p\) drift norms
then gives the increment estimate. The two summands of \(c\) in
\eqref{eq:factor} have total variations \(\abs\alpha T\) and
\(\beta^2T^{2H}/2\). In particular the terminal singularity of \(c'\)
when \(H<1/2\) causes no loss after integration.
\end{proof}

\section{Conditional quadrature for fractional noise}\label{sec:fractional}

Throughout this section \(H\ne1/2\). We realize the normalized noise using
a two-sided Brownian motion \(W\):
\begin{equation}\label{eq:MVN}
 B^H_t=\gamma_H\int_{\R}
       \bigl((t-u)_+^{H-1/2}-(-u)_+^{H-1/2}\bigr)\,dW_u ,
\end{equation}
where \(\gamma_H>0\) is chosen so that \(\E(B^H_1)^2=1\).
The integrand is square integrable: the finite singularities have squared
exponent \(2H-1>-1\), and its difference at negative infinity is
\(O(\abs u^{H-3/2})\). The Wiener integral and its isometry are as in
\cite[Theorem 3.7.1]{Chafai}; the negative and positive half-lines use
independent Brownian motions. Translation gives stationary increments and
dilation gives variance \(t^{2H}\), so polarization yields exactly the
covariance in Section~\ref{sec:setup}. Gaussian moment bounds and the
continuity criterion of \cite[Theorem 3.6.3]{Chafai} give a continuous
version. Since the exact terminal value and its approximation
in \eqref{eq:exact-ode}--\eqref{eq:scheme} are functions of the same
continuous Gaussian path, this realization preserves their joint law.

Let \(\F_s=\sigma(W_u:u\le s)\), completed, and write
\(\E^s=\E[\cdot\mid\F_s]\).
This is the driving Brownian filtration, including the negative-time past;
we do not identify it with the natural filtration of \(B^H\).
For \(r>s\), the conditional Gaussian mean and variance are
\[
 m_s(r)=\E^s B^H_r,\qquad
 v_s(r)=q_H(r-s)^{2H},\qquad q_H=\frac{\gamma_H^2}{2H}.
\]
The innovation is independent of \(\F_s\). Thus, for an
\(\F_s\)-measurable \(z\),
\begin{equation}\label{eq:conditional-heat}
 \E^s F(c,B^H_r,z)=P_{v_s(r)}f_{c,z}(m_s(r)).
\end{equation}
For \(s\le u<r\) with \(r-u\le u-s\), the Wiener isometry and the
mean-value formula for the kernel give
\begin{equation}\label{eq:conditional-mean}
 \norm{m_s(r)-m_s(u)}_p
 \le C_p(r-u)(u-s)^{H-1}.
\end{equation}
Indeed, the squared kernel derivative integrates on \((-\infty,s]\)
to a constant times \((u-s)^{2H-2}\); Gaussian scaling gives every
finite \(p\). Moreover \(\operatorname{Var}(m_s(r))\le T^{2H}\).
Its exponential moments are therefore uniformly bounded.

The slow Euler state can be predicted more accurately than the noise.

\begin{lemma}\label{lem:prediction}
For \(p\ge1\) and \(0\le s<t\le T\),
\begin{equation}\label{eq:prediction}
 \norm{Z^n_{\eta(t)}-\E^s Z^n_{\eta(t)}}_p
 \le C_p(t-s)^{1+H}.
\end{equation}
\end{lemma}
\begin{proof}
The left side is zero if \(\eta(t)\le s\).
Otherwise put \(s^+=h\lceil s/h\rceil\le\eta(t)\).
The state \(Z^n_{s^+}\) is \(\F_s\)-measurable: the full Euler step
ending at \(s^+\) uses the preceding grid noise and state.
The random variable
\[
 V=Z^n_{s^+}+
       \int_{s^+}^{\eta(t)}F(c(\eta(r)),B^H_s,Z^n_{s^+})\,dr
\]
is \(\F_s\)-measurable. Keeping the deterministic \(c(\eta(r))\)
unchanged, \eqref{eq:basic-F}, \eqref{eq:F-noise} and
Lemma~\ref{lem:moments} yield
\[
 \norm{Z^n_{\eta(t)}-V}_p
 \le C_p\int_{s^+}^{\eta(t)}
       \bigl(\eta(r)-s^++(\eta(r)-s)^H\bigr)\,dr
 \le C_p(t-s)^{1+H}.
\]
Conditional contraction proves \eqref{eq:prediction}, with a factor two.
In particular, if \(s_1\le s_2\le r\) and \(r-s_1\le C_0\ell\), then
\begin{equation}\label{eq:prediction-difference}
 \norm{\E^{s_1}Z^n_{\eta(r)}-\E^{s_2}Z^n_{\eta(r)}}_p
 \le C_p\ell^{1+H}.
\end{equation}
This follows by conditioning \(Z^n_{\eta(r)}-\E^{s_1}Z^n_{\eta(r)}\)
at \(s_2\).
\end{proof}

We estimate a signed quadrature, rather than the sum of absolute cell
errors. Put \(c_r=c(\eta(r))\), \(U_r=Z^n_{\eta(r)}\), and
\begin{equation}\label{eq:GQ}
 G_r(z)=F(c_r,B^H_r,z)-F(c_r,B^H_{\eta(r)},z),\qquad
 Q_h(t)=\int_0^tG_r(U_r)\,dr.
\end{equation}

\begin{proposition}\label{prop:fractional-Q}
Uniformly for \(0\le t\le T\),
\begin{equation}\label{eq:fractional-Q}
 \norm{Q_h(t)}_4\le
 \begin{cases}
 C h^{H+1/2},&0<H<1/2,\\
 C h,&1/2<H<1 .
 \end{cases}
\end{equation}
\end{proposition}
\begin{proof}
For an interval \([s,t]\) of length \(\ell=t-s\), let
\(v=(s-\ell)\vee0\) and define the conditional increment
\begin{equation}\label{eq:A}
 \Acal_{s,t}=\E^v\int_s^tG_r(\E^v U_r)\,dr.
\end{equation}
We use intervals with \(s=0\) or \(s\ge\ell\), which include every
dyadic interval in \([0,t]\). For every fixed \(p\ge1\),
\begin{equation}\label{eq:A-size}
 \norm{\Acal_{s,t}}_p
 \le C_p b_h(\ell),\qquad
 b_h(\ell)=\min\{\ell h^H,h\ell^H\}.
\end{equation}
The first bound is \eqref{eq:F-noise} integrated over the interval.
For the second, suppose first that \(s\ge\ell\) and \(\ell\ge2h\).
Then \(\eta(r)-v\ge\ell/2\). Conditional Gaussian convolution,
\eqref{eq:conditional-mean}, and Lemma~\ref{lem:heat} give, with
\(z_r=\E^v U_r\),
\[
 \begin{split}
 &\norm{P_{v_v(r)}f_{c_r,z_r}(m_v(r))
       -P_{v_v(\eta(r))}f_{c_r,z_r}(m_v(\eta(r)))}_p\\
 &\hspace{25mm}\le
 C_p\bigl(h\ell^{H-1}+\ell^{-H}h\ell^{2H-1}\bigr)
 \le C_ph\ell^{H-1}.
 \end{split}
\]
The two terms respectively change the mean and variance.
Their weighted derivative factors have bounded higher moments, by
Lemma~\ref{lem:moments} and the Gaussian bounds on \(m_v\).
Integrating gives \(C_ph\ell^H\). For \(s=0\), the part \(r\le2h\)
costs \(C_ph^{1+H}\); on \(r>2h\) the same integrand estimate is
\(C_ph r^{H-1}\), whose integral is \(C_ph\ell^H\).
If \(h\le\ell\le2h\), the first bound implies the second up to a
fixed factor; if \(\ell<h\), the minimum is the first bound.

Let \(u=(s+t)/2\) and
\(\delta\Acal_{s,u,t}=\Acal_{s,t}-\Acal_{s,u}-\Acal_{u,t}\).
For \(s\ge\ell\) we also have
\begin{equation}\label{eq:A-delta}
 \norm{\E^{s-\ell}\delta\Acal_{s,u,t}}_p
 \le C_ph\ell^{1+H}.
\end{equation}
If \(\ell\le4h\), the tower property expresses the conditional
difference as integrals of \(G_r(z_1)-G_r(z_2)\), where \(z_1\)
and \(z_2\) are the parent and child predictions of \(U_r\).
The \(2L\)-Lipschitz bound on \(G_r\) and
\eqref{eq:prediction-difference} give
\(C_p\ell^{2+H}\le C_ph\ell^{1+H}\).

For \(\ell\ge4h\), the child's conditioning time is
\(w=s-\ell/2\) on the left half and \(w=s\) on the right half.
On its corresponding half, \(\eta(r)-w\ge\ell/4\).
Both \(z_1=\E^{s-\ell}U_r\) and \(z_2=\E^wU_r\) are
\(\F_w\)-measurable. Apply \eqref{eq:conditional-heat} to the bounded
function \(g=f_{c_r,z_1}-f_{c_r,z_2}\).
Equations \eqref{eq:heat-mixed} and \eqref{eq:conditional-mean}
bound the conditional integrand in \(L^p\) by
\[
 C_p\norm{z_1-z_2}_{2p}
       \bigl(\ell^{-H}h\ell^{H-1}
                    +\ell^{-2H}h\ell^{2H-1}\bigr)
 \le C_ph\ell^H.
\]
Integrate over the two halves and condition once more at \(s-\ell\)
to obtain \eqref{eq:A-delta}. If \(s=0\), the parent and both children
condition at zero, so \(\delta\Acal_{0,t/2,t}=0\) exactly.

For the dyadic subdivision of \([0,t]\) at length
\(\ell_j=t2^{-j}\), let \(S_j\) be the sum of its conditional
increments. The difference \(S_{j+1}-S_j\) is minus the sum of their
\(\delta\Acal\)'s. At the interval \([i\ell_j,(i+1)\ell_j]\),
this difference is \(\F_{i\ell_j}\)-measurable, since the latest
child conditioning time is \(i\ell_j\). For \(i\ge1\), subtract its
conditional expectation at \(\F_{(i-1)\ell_j}\).
The resulting sequence is a martingale-difference sequence; the
initial interval contributes zero. Equations
\eqref{eq:A-size}--\eqref{eq:A-delta} give
\begin{equation}\label{eq:scale-sum}
 \norm{S_{j+1}-S_j}_4
 \le C\sqrt{t/\ell_j}\,b_h(\ell_j)+Cth\ell_j^H.
\end{equation}
For completeness, the required discrete estimate is elementary.
If \(D_i\) are martingale differences, \(M=\sum_i\norm{D_i}_4^2\),
and \(S=\sum_iD_i\), then
\[
 S^2=2\sum_iS_{i-1}D_i+\sum_iD_i^2,\qquad
 \norm S_4^2\le2\norm S_4\sqrt M+M.
\]
The first sum has orthogonal \(L^2\) increments, and conditional
Jensen bounds every \(\norm{S_{i-1}}_4\) by \(\norm S_4\).
Solving the quadratic gives
\(\norm S_4\le(1+\sqrt2)\sqrt M\), which proves
\eqref{eq:scale-sum}.

The sums \(S_j\) converge to \(Q_h(t)\) in \(L^4\).
On an interval of length \(\ell\), replacing \(U_r\) by its prediction
costs \(C\ell^{2+H}\), by \eqref{eq:prediction}.
Conditional projection of \(F(c_r,B^H_r,z_r)\) costs \(C\ell^H\)
pointwise: its innovation has variance at most \(C\ell^{2H}\),
and the noise increment bound has integrable exponential weights.
The same holds for \(B^H_{\eta(r)}\), with zero cost if
\(\eta(r)\le v\). Hence
\[
 \norm{S_j-Q_h(t)}_4\le CT(\ell_j^H+\ell_j^{1+H})\longrightarrow0.
\]

In \eqref{eq:scale-sum}, scales \(\ell_j\ge h\) contribute
\[
 C\sqrt t\,h\sum_{\ell_j\ge h}\ell_j^{H-1/2},
\]
whereas scales \(\ell_j<h\) contribute at most
\(C\sqrt t\,h^{H+1/2}\).
The conditional-mean terms sum to
\(Cth\sum_j\ell_j^H\le Ch\).
The initial increment satisfies \eqref{eq:A-size}.
The first geometric sum is dominated by its smallest scale if
\(H<1/2\) and by its largest scale if \(H>1/2\), giving
\eqref{eq:fractional-Q}. For \(t<h\), the direct estimate
\(Cth^H\) suffices. All bounds are uniform for \(t\le T\).
\end{proof}

At \(H=1/2\), the first geometric sum is critical and this argument
would give \(h(1+\log(T/h))\). We next use the Brownian structure to
recover order one without differentiating the drift.

\section{The Brownian endpoint}\label{sec:brownian}

Here \(H=1/2\), \(B^{1/2}=W\) is ordinary Brownian motion, and
\(c(t)=(-\alpha+\beta^2/2)t\).
We first record the quadrature estimate that replaces the critical scale
sum. The stochastic-calculus facts used in this section are ordinary
It\^o calculus, the Brownian martingale characterization, Girsanov's
theorem for bounded progressive drifts, and the
Burkholder--Davis--Gundy inequalities; see
\cite[Theorems 7.1.1, 7.2.1, 7.5.1 and 7.7.1]{Chafai}.

\begin{lemma}\label{lem:Brownian-quadrature}
Let \(W_0=0\), \(d=mh\) for an integer \(m\ge1\), and \(f:\R\to\R\)
be globally Lipschitz. For every \(p\ge2\),
\begin{equation}\label{eq:Brownian-quadrature}
 \left\|\sup_{t\le d}\left|
   \int_0^t(f(W_s)-f(W_{\eta(s)}))\,ds\right|\right\|_p
 \le C_p[f]_{\mathrm{Lip}}h\sqrt d .
\end{equation}
The bound is independent of \(f(0)\) and of shifts of its argument.
\end{lemma}
\begin{proof}
First suppose \(f\) is smooth with bounded first derivative.
Set \(g_h(s)=(\lfloor s/h\rfloor+1)h-s\).
At grid endpoints the quadrature is \(\int_0^t g_h(s)\,df(W_s)\).
At intermediate \(t\), its difference from that integral is
\[
 -\bigl((\lfloor t/h\rfloor+1)h-t\bigr)
          (f(W_t)-f(W_{\eta(t)})).
\]
The maximum of this correction has \(L^p\) norm at most
\(C_p[f]_{\mathrm{Lip}}h\sqrt d\).

The change \(df(W)\) is the average of its forward and right-endpoint
backward stochastic differentials:
\begin{equation}\label{eq:forward-backward}
 \int_0^t g_h\,df(W)
 =\frac12\int_0^t g_h(s)f'(W_s)\,dW_s
  +\frac12\int_0^t g_h(s)f'(W_s)\,d^-W_s.
\end{equation}
The forward integral equals the change in \(f(W)\) minus
one-half the \(f''\) integral, while the backward integral equals that
change plus one-half the \(f''\) integral. The deterministic
bounded-variation function \(g_h\) adds no quadratic covariation.
This proves \eqref{eq:forward-backward} first on each cell and then
on their union.

To estimate the backward integral, let \(Y_u=W_{d-u}\) in its natural
reverse filtration. Its Gaussian bridge transition gives, for \(u<v<d\),
\[
 \E[Y_v\mid Y_r,\ r\le u]=\frac{d-v}{d-u}Y_u.
\]
It follows that
\(M_u=Y_u-Y_0+\int_0^uY_r/(d-r)\,dr\) is a continuous martingale.
Its quadratic variation is \(u\), because the added integral has
finite variation and Brownian increments retain their quadratic
variation on reversal. The Brownian martingale characterization
therefore gives
\[
 dY_u=dM_u-\frac{Y_u}{d-u}\,du.
\]
The drift extends to the endpoint in every finite \(L^p\), since
\(\int_0^d\norm{Y_u}_p/(d-u)\,du=2C_p\sqrt d\).
Writing
\(N_u=\int_0^u g_h(d-v)f'(Y_v)\,dM_v\), the reversal is explicitly
\[
 \int_0^t g_h(s)f'(W_s)\,d^-W_s
 =-\int_{d-t}^d g_h(d-u)f'(Y_u)\,dY_u
 =-(N_d-N_{d-t})
   +\int_{d-t}^d g_h(d-u)f'(Y_u)\frac{Y_u}{d-u}\,du .
\]
The martingale increment is controlled by \(2\sup_{u\le d}|N_u|\);
its integrand is bounded by \(h[f]_{\mathrm{Lip}}\).
The drift term is bounded by
\[
 h[f]_{\mathrm{Lip}}\int_0^d\frac{\abs{W_s}}s\,ds,
\]
whose \(L^p\) norm is at most \(2C_p h[f]_{\mathrm{Lip}}\sqrt d\).
The forward martingale maximum has the same bound by the maximal
martingale inequality. One can first reverse before \(d\) and then
pass to \(d\) using the displayed integrability bounds.
This proves \eqref{eq:Brownian-quadrature}.
Finally, mollify \(f\) by a Gaussian kernel. The Lipschitz constant
does not increase and the approximations converge uniformly, so their
quadratures converge uniformly. The same estimate holds for \(f\).
\end{proof}

\begin{lemma}\label{lem:bounded-drift}
Suppose \(\beta\ne0\) and, on a grid block of length \(d=mh\),
\[
 dq_t=u(q_t)\,dt+\beta\,dW_t,\qquad \abs u\le M.
\]
If \(M^2d\le C_0\), then for every bounded Lipschitz \(b\) and \(p\ge2\),
\begin{equation}\label{eq:bounded-drift}
 \left\|\sup_{t\le d}\left|
   \int_0^t(b(q_s)-b(q_{\eta(s)}))\,ds\right|\right\|_p
 \le C_{p,C_0,\beta}[b]_{\mathrm{Lip}}h\sqrt d .
\end{equation}
The constant is uniform in the initial value. The same estimate holds
conditionally at a stopping time in a finite deterministic grid, with
the initial value and block length known at that time.
\end{lemma}
\begin{proof}
Let \(\mu_t=u(q_t)/\beta\) and
\(D=\exp(-\int_0^d\mu_t\,dW_t-\tfrac12\int_0^d\mu_t^2\,dt)\).
The bounded progressive coefficient makes this a probability density.
Under \(d\mathbb Q=D\,d\mathbb P\), the process
\(\widetilde W_t=W_t+\int_0^t\mu_s\,ds\) is Brownian and
\(q_t=q_0+\beta\widetilde W_t\).
For each fixed moment exponent, both the forward and inverse densities
have moments bounded by \(\exp(C_pM^2d/\beta^2)\); this follows by
factoring their powers into an exponential martingale and its
bounded quadratic-variation correction.
H\"older's inequality and Lemma~\ref{lem:Brownian-quadrature},
applied at exponent \(2p\) to \(f(x)=b(q_0+\beta x)\), prove
\eqref{eq:bounded-drift}.
At a finite-grid stopping time, condition on each of its possible
values. The future Brownian increments are independent of the past,
and all block parameters are already known. The same density argument
gives the conditional assertion.
\end{proof}

We retain the original \(Z\)-equation for stability.
The logarithm will only be used on intervals where \(Z\) stays away
from zero. On a sign interval \(\sigma=\operatorname{sign}Z\), recall
\[
 b_\sigma(q)=\frac{a(\sigma e^q)-a_0}{\sigma e^q},
 \qquad Z'_t=J_ta(Z_t/J_t).
\]
For \(q_t=\log\abs{Z_t/J_t}=\log\abs{Z_t}-c(t)+\beta W_t\),
the finite-variation term has differential \((Z'_t/Z_t)\,dt\).
Since \(c'=-\alpha+\beta^2/2\), this gives
\begin{equation}\label{eq:log-SDE}
 dq_t=\left(\alpha-\frac{\beta^2}{2}
          +b_\sigma(q_t)+\sigma a_0e^{-q_t}\right)\,dt
          +\beta\,dW_t.
\end{equation}
If \(a_0=0\) and \(x_0\ne0\), then
\(\abs{x_0}e^{-LT}\le\abs{Z_t}\le\abs{x_0}e^{LT}\), so its sign
never changes. The drift in \eqref{eq:log-SDE} is bounded.
The noise quadrature vanishes when \(\beta=0\), so assume \(\beta\ne0\) here.
Define
\[
 R(t)=\int_0^t[b_\sigma(q_s)-b_\sigma(q_{\eta(s)})]\,ds .
\]
Lemma~\ref{lem:bounded-drift} over the whole interval gives
\(\|\sup_{t\le T}|R(t)|\|_p\le C_ph\).
By \eqref{eq:basic-F}, \eqref{eq:F1-bounds} and
Lemma~\ref{lem:moments}, the exact-state \(F_1\) noise quadrature on a
grid prefix, defined in \eqref{eq:Brownian-Q} below, differs from
\(\int_0^{t_k}Z_{\eta(s)}\,dR(s)\) by \(O(h)\) in \(L^p\).
Summation by parts bounds this weighted integral by
\[
 \bigl(\sup_{t\le T}|Z_t|+\TV(Z_\eta;[0,T])\bigr)
       \sup_{t\le T}|R(t)|.
\]
The first factor is deterministically bounded, so the \(F_1\) noise
quadrature has order \(h\). If also \(x_0=0\), the solution and scheme
vanish identically.

When \(a_0\ne0\), the last drift term in \eqref{eq:log-SDE} can be
singular. Short blocks chosen by the distance of \(Z\) from zero
keep its density cost bounded. The resulting estimates grow only
polynomially in the integrating factor.

\begin{proposition}\label{prop:Brownian-Q}
For \(H=1/2\), uniformly at grid times \(t_k\),
\begin{equation}\label{eq:Brownian-Q}
 \left\|\int_0^{t_k}
 \bigl(F_1(c(\eta(s)),W_s,Z_{\eta(s)})
       -F_1(c(\eta(s)),W_{\eta(s)},Z_{\eta(s)})\bigr)\,ds
 \right\|_4\le Ch .
\end{equation}
\end{proposition}
\begin{proof}
The case \(a_0=0\) was just discussed, with the \(O(h)\) state and
\(c\) freezing bounds from \eqref{eq:F1-bounds} and
Lemma~\ref{lem:moments}. If \(\beta=0\), the integral vanishes.
Reflection \(X\mapsto-X\) permits us to assume \(a_0>0\).
If \(L=0\), then \(a_1=0\), so suppose \(L>0\).
Let
\[
 M_J=\max\{1,\sup_{t\le T}J_t,\sup_{t\le T}J_t^{-1}\},
 \qquad E_K=\{M_J\le K\},\quad K\ge2.
\]
On \(E_K\), deterministic Gronwall gives
\begin{equation}\label{eq:localized-Z}
 \sup_{t\le T}\abs{Z_t}\le C_ZK,\qquad
 \sup_{t\le T}\abs{Z'_t}\le C_DK.
\end{equation}
Choose a fixed \(r_*>0\) with
\(r_*\le\min\{1,a_0/(4L)\}\) and put \(r_0=r_*/K\).
In the band \(\abs{Z_t}\le2r_0\),
\begin{equation}\label{eq:monotone-band}
 Z'_t=a_0J_t+F_1(c(t),W_t,Z_t)\ge\frac{a_0}{2K}.
\end{equation}
Thus \(Z\) traverses the band monotonically. A positive or negative
annulus inside the band is traversed at most once; the time spent
in \(\epsilon/2\le\abs Z\le3\epsilon\) is at most \(CK\epsilon\)
when \(3\epsilon\le2r_0\).

Choose a sufficiently large fixed \(D\ge8C_DC_Z+1\), and put
\(\delta=2\sqrt D\,K\sqrt h\).
We first assume \(h\le c_*K^{-4}\), where the fixed \(c_*>0\) is
small enough that \(2\delta\le r_0\) and \(C_DKh\le\delta\).
Partition a prefix \([0,t_k]\) at grid times as follows.
A start \(t_i\) with \(\abs{Z_{t_i}}\le\delta\) receives one cell.
Otherwise choose \(\epsilon=\delta2^j\) with
\(\epsilon\le\abs{Z_{t_i}}<2\epsilon\), and take
\begin{equation}\label{eq:block-length}
 m_i=\min\left\{k-i,\left\lfloor\frac{\epsilon^2}{DK^2h}
                         \right\rfloor\right\}
\end{equation}
consecutive cells. Every start is a finite-grid stopping time.
Its sign, \(\epsilon\), and length \(d=m_ih\) are known at that time.
The floor is at least four before the final truncation, and
\[
 d\le\frac{\epsilon^2}{DK^2},\qquad
 d\ge\frac{\epsilon^2}{2DK^2}
\]
unless this is the last block.

On \(E_K\), \(\epsilon\le C_ZK\), and \eqref{eq:localized-Z} gives
\(\abs{Z_t-Z_{t_i}}\le C_DKd\le\epsilon/4\).
Throughout an outer block, \(Z\) therefore keeps its sign and
\(\epsilon/2\le\abs Z\le3\epsilon\), while
\(\abs{Z/J}\ge\epsilon/(2K)\).
Clip the last term of \eqref{eq:log-SDE} to
\[
 \sigma\min\{a_0e^{-q},2a_0K/\epsilon\}.
\]
This produces a globally Lipschitz bounded drift agreeing with the
exact \(q\) on the block on \(E_K\).
For starts with \(\epsilon\le C_ZK\), its bound is \(CK/\epsilon\),
and its squared bound times \(d\) is uniformly bounded.
Let \(q^{\mathrm{cl}}\) be this clipped process with
\(q^{\mathrm{cl}}_{t_i}=q_{t_i}\). For this block \(b=[t_i,t_i+d]\), put
\[
 R_b^{\mathrm{cl}}(u)
   =\int_{t_i}^u
       [b_\sigma(q^{\mathrm{cl}}_s)
        -b_\sigma(q^{\mathrm{cl}}_{\eta(s)})]\,ds,
 \qquad t_i\le u\le t_i+d .
\]
On \(E_K\), write \(R_b\) for the same primitive using the unclipped \(q\);
the supremum below is over this block. Lemma~\ref{lem:bounded-drift},
with test function \(b_\sigma\), gives conditionally at the start
\begin{equation}\label{eq:clipped-R}
 \left(\E\left[\sup\abs{R_b^{\mathrm{cl}}}^p
                    \mid\F_{t_i}\right]\right)^{1/p}
 \le C_ph\sqrt d\le C_ph\epsilon/K .
\end{equation}
The Lipschitz constant of \(b_\sigma\) is at most \(2L\), independently
of the clipping level.

On this block, the integrand in \eqref{eq:Brownian-Q} differs from
\(Z_{\eta(s)}[b_\sigma(q_s)-b_\sigma(q_{\eta(s)})]\) by at most
\(CKh\) on \(E_K\). Indeed, \eqref{eq:F1-bounds} and
\eqref{eq:localized-Z} freeze \(Z_s\) and the linear \(c(s)\)
at cost \(CKh\). These errors over all blocks cost \(CKh\).
Summation by parts bounds the weighted quadrature by
\[
 \bigl(\sup\abs Z+\TV(Z_\eta;\text{block})\bigr)\sup\abs{R_b}
 \le C\epsilon\,\sup\abs{R_b},
\]
since \(\TV(Z_\eta;\text{block})\le C_DKd\le\epsilon/4\).
On \(E_K\) the original and clipped processes agree.
Thus each block contribution, multiplied by \(\ind_{E_K}\),
has \(L^p\) norm at most \(C_ph\epsilon^2/K\).


For \(\epsilon<r_0/4\), every such block lies in its expanded annulus.
Equations \eqref{eq:monotone-band} and \eqref{eq:block-length}
bound the number of these disjoint blocks by \(CK^3/\epsilon\),
apart from one final truncated block.
Their total \(L^p\) contribution is bounded by
\[
 C_phK^2\epsilon+C_ph\epsilon^2/K .
\]
Summing the dyadic \(\epsilon\)'s below \(r_0/4\) gives a polynomial
multiple of \(h\). For the remaining outer blocks,
\(\epsilon\ge r_0/4\); every nonfinal length is at least \(cK^{-4}\).
There are at most \(CK^4\) of them, and each contribution is at most
\(C_phK\) because \(\epsilon\le C_ZK\) on \(E_K\).
They therefore cost at most \(C_phK^5\).

These counts remain valid in the norm estimates even though the block
starts are random. For a fixed small annulus, enumerate its successive
starts as grid stopping times. On \(E_K\), only the first deterministic
upper bound on their number is present. Each contribution with its
indicator is bounded by \(C\epsilon\sup\abs{R_b^{\mathrm{cl}}}\);
\eqref{eq:clipped-R} gives its norm conditionally at that start.
For large blocks enumerate only starts with \(\epsilon\le C_ZK\),
a condition already known at the start. On \(E_K\) this includes
every large block. Apply the deterministic count bound and
Minkowski's inequality. Missing starts contribute zero at \(t_k\).
No independence of \(E_K\) and the block errors is used.

It remains to bound the cells starting within \(\abs{Z_{\eta(s)}}\le\delta\).
On \(E_K\), all their times lie within \(\abs Z\le2\delta\), by
\(C_DKh\le\delta\). Their total length is at most \(CK\delta\).
The noise increment estimate \eqref{eq:F1-bounds} bounds their
quadrature pathwise by
\[
 C\delta(K\delta)D_h\le CK^3hD_h,\qquad
 D_h=\max_{i<n}\sup_{s\in[t_i,t_{i+1}]}\abs{W_s-W_{t_i}} .
\]
The Brownian maximal tail bound and a union bound give
\(\norm{D_h}_p\le C_p\sqrt{h\log(2n)}\), uniformly bounded as
\(h=T/n\). This inner contribution is consequently at most
\(C_phK^3\) in \(L^p\).
Writing \(Q_h^{(1)}(t_k)\) for the integral in
\eqref{eq:Brownian-Q}, we have proved
\begin{equation}\label{eq:polynomial-localization}
 \norm{Q_h^{(1)}(t_k)\ind_{E_K}}_p\le C_phK^6,\qquad p\ge2 .
\end{equation}
The harmless enlargement from \(K^5\) to \(K^6\) includes all preceding
constants. If \(h>c_*K^{-4}\), the direct increment and moment bounds
give \(\norm{Q_h^{(1)}(t_k)}_p\le C_p\sqrt h\le C_phK^2\),
so \eqref{eq:polynomial-localization} holds for every \(n\).

Finally \(J_t=e^{c(t)-\beta W_t}\), and the Brownian supremum tail implies
\[
 \mathbb P(M_J>K)\le C\exp\{-c_1(\log K-C_1)^2\}
\]
for large \(K\), with fixed constants.
For dyadic \(K\), H\"older's inequality and
\eqref{eq:polynomial-localization} at \(p=8\) give
\[
 \norm{Q_h^{(1)}(t_k)
        \ind_{\{K/2<M_J\le K\}}}_4
 \le ChK^6\mathbb P(M_J>K/2)^{1/8}.
\]
The dyadic sum is finite, independently of \(n\), proving
\eqref{eq:Brownian-Q}.
\end{proof}

The constant drift part is \(a_0J_t\).
Ordinary It\^o calculus gives
\begin{equation}\label{eq:dJ}
 dJ_t=(-\alpha+\beta^2)J_t\,dt-\beta J_t\,dW_t.
\end{equation}
Integrating this differential with the deterministic sawtooth weight
bounded by \(h\), and using Gaussian exponential moments and the
maximal martingale inequality, shows
\begin{equation}\label{eq:J-quadrature}
 \left\|\int_0^{t_k}(J_s-J_{\eta(s)})\,ds\right\|_4\le Ch .
\end{equation}
Combining \eqref{eq:Brownian-Q} and \eqref{eq:J-quadrature} therefore
controls the entire Brownian drift quadrature.

\begin{proof}[Proof of Theorem~\ref{thm:main}]
Let \(e_k=Z_{t_k}-Z^n_{t_k}\) and
\(\rho=\min\{H+1/2,1\}\).
For \(H\ne1/2\), subtract \eqref{eq:interpolation} from
\eqref{eq:exact-ode} and use the decomposition
\[
 \begin{split}
 F(c(s),B^H_s,Z_s)-F(c(\eta(s)),B^H_{\eta(s)},Z^n_{\eta(s)})
  ={}&F(c(s),B^H_s,Z_s)-F(c(s),B^H_s,Z^n_{\eta(s)})\\
   &+F(c(s),B^H_s,Z^n_{\eta(s)})
                 -F(c(\eta(s)),B^H_s,Z^n_{\eta(s)})\\
   &+G_s(Z^n_{\eta(s)}).
 \end{split}
\]
The first line costs \(Ch+Lh\sum_{i<k}\norm{e_i}_4\), by
\eqref{eq:basic-F} and Lemma~\ref{lem:moments}.
The \(c\) increments cost \(Ch\TV(c)\): their \(L^4\) weighted
factors are uniformly bounded, and
\[
 \int_0^{t_k}\abs{c(s)-c(\eta(s))}\,ds\le h\TV(c).
\]
Proposition~\ref{prop:fractional-Q} bounds the last line by \(Ch^\rho\).
Hence
\begin{equation}\label{eq:Gronwall-error}
 \norm{e_k}_4\le Ch^\rho+Lh\sum_{i<k}\norm{e_i}_4 .
\end{equation}

For \(H=1/2\), use instead the exact-state residual
\[
 R_k=\int_0^{t_k}
  \bigl(F(c(s),W_s,Z_s)-F(c(\eta(s)),W_{\eta(s)},Z_{\eta(s)})\bigr)\,ds.
\]
Freezing \(Z_s\) and the linear \(c(s)\) costs \(Ch\).
The remaining \(a_1\) quadrature is
Proposition~\ref{prop:Brownian-Q}; the \(a_0\) quadrature is
\eqref{eq:J-quadrature}. Thus \(\norm{R_k}_4\le Ch\).
The deterministic \(z\)-Lipschitz estimate at each grid point gives
\eqref{eq:Gronwall-error} with \(\rho=1\).
The case \(\beta=0\) follows directly from deterministic Euler
stability and consistency.

Discrete Gronwall yields \(\norm{e_n}_4\le Ch^\rho\).
The prescribed quotient uses the same strictly positive \(J_T\)
for the exact and approximate terminal states. Therefore
\[
 \norm{X_T-X_n}_2\le\norm{e_n}_4\norm{J_T^{-1}}_4\le Ch^\rho,
\]
as required. These estimates may first be used for \(h\le\min\{T,1\}\).
The finitely many larger steps have finite errors by the moment bounds,
and enlarging the fixed-data constant includes them.
If \(a=0\), both transformed states are identically \(x_0\) and
the scheme is exact.
\end{proof}

\paragraph{AI use.}
OpenAI Codex was used for proof development, source checking and drafting
this manuscript. Independent mathematical review is requested.

\begin{thebibliography}{9}
\bibitem{ELS}
U.~Erdo\u{g}an, G.~J. Lord and R.~B. Schieven,
\emph{Numerical approximation of SDEs driven by fractional Brownian motion
for all \(H\in(0,1)\) using WIS integration},
IMA Journal of Numerical Analysis (2026),
\href{https://doi.org/10.1093/imanum/drag055}{doi:10.1093/imanum/drag055}.
Version used: \href{https://arxiv.org/abs/2404.07013v2}{arXiv:2404.07013v2}.

\bibitem{BDG}
O.~Butkovsky, K.~Dareiotis and M.~Gerencs\'er,
\emph{Approximation of SDEs: a stochastic sewing approach},
Probability Theory and Related Fields \textbf{181} (2021), 975--1034,
\href{https://doi.org/10.1007/s00440-021-01080-2}{doi:10.1007/s00440-021-01080-2}.

\bibitem{Moritoki}
T.~Moritoki,
\emph{Strong rate of convergence for the Euler--Maruyama scheme of additive
fractional SDEs with Lipschitz drift},
\href{https://arxiv.org/abs/2607.24166v1}{arXiv:2607.24166v1} (2026).

\bibitem{Altmeyer}
R.~Altmeyer,
\emph{Approximation of occupation time functionals},
\href{https://arxiv.org/abs/1706.03418v4}{arXiv:1706.03418v4} (2021).

\bibitem{KR}
P.~E. Kloeden and R.~M.~S. Rosa,
\emph{Strong order-one convergence of the Euler method for random ordinary
differential equations driven by semi-martingale noises},
\href{https://arxiv.org/abs/2306.15418v8}{arXiv:2306.15418v8} (2025).

\bibitem{Chafai}
D.~Chafa\"i,
\emph{Introduction to stochastic calculus},
Universit\'e Paris-Dauphine--PSL, 2020--2021 lecture notes,
corrected version compiled December 7, 2025,
\url{https://djalil.chafai.net/docs/M2/m2-stochastic-calculus-course-2020-2021.pdf}.
\end{thebibliography}
\end{document}
